%%%%%%%%%%%%%%%%CONDENSED REWRITE

\section{\toolname{} Design}
\label{sec:design}

We first describe how \toolname{} predicts $\hat{y}_t$, then (Section~\ref{sec:gates}) the gates that determine the auditing policy $\pi$.

\toolname  is designed around two core empirical observations. First, although a task's true data consumption $c_t$ cannot be observed by the scheduler before launch, it can be estimated using available prelaunch features and memory metrics. This lightweight, learned proxy $\hat{c}_t$ serves as a critical supporting signal for peak memory estimation.

Second, accurate memory allocation requires learning at the correct level of task specificity. Frequently repeated pipeline steps often possess dozens of nearly identical historical runs to learn from, whereas a newly deployed workflow stage may have none. Monolithic global models degrade under this variance by averaging over these non-uniform historical densities.

Central to the design of \toolname{} is a hierarchical system of \textit{scopes}. A \textit{scope} is a task compatibility criterion defined by task instance, staged input, structural features, workflow name, process type, software version, and underlying hardware profile. Scopes are used to isolate pools of prior completed tasks whose peak-memory and consumption histories are relevant to a pending execution. \toolname{} defines $K$ distinct scopes, $I_0, \dots, I_{K-1}$, ordered strictly from most specific (e.g., exact task instance previously executed) to most general (e.g., matching cluster hardware profile). A task uses exactly one active scope based on historical support, or a cold-start state when no scope is supported, but its feature space remains enriched by historical summaries across all scopes.

\subsection{The Predictor Pipeline}

The problem formulation seeks a single allocation function
$f(\mathbf{s}_t, \mathbf{h}_t)$. \toolname{} instantiates $f$ as a prediction
pipeline: a global point model, process-specific upper-tail models with broader
fallbacks, an auxiliary consumption estimator $g$, and a tail calibration
layer $\mathcal{K}$. The pipeline executes sequentially:
\begin{equation}
\hat{c}_t = g(\mathbf{s}_t, \mathbf{h}_{m,t})
\label{eq:proxy_consumption}
\end{equation}
\begin{equation}
\hat{y}_t = \mathcal{K}\Big( f\big(\mathbf{s}_t,\,
\mathbf{h}_t,\, \hat{c}_t\big),\,\mathbf{s}_t,\,\mathbf{h}_t \Big).
\label{eq:final_allocation}
\end{equation}
The proxy $\hat{c}_t$ is computed from prelaunch-legal inputs
$(\mathbf{s}_t, \mathbf{h}_{m,t})$ alone, so the pipeline uses no
information beyond the inputs admitted by the problem statement. The
asymmetric allocation objective $L$ is implemented through quantile pinball
loss and allocation-policy selection. The prediction models estimate the
point and upper-tail quantiles of $y_t$, while the calibration layer
$\mathcal{K}$ (Section~\ref{sec:calibration}) selects a base request and
corrects residual under-coverage.

\subsection{Scope Hierarchy and Routing}

Let $\mathcal{P}(t)$ denote the causal universe of all tasks that have fully completed prior to the launch epoch of task t. 

For each hierarchy tier $i \in \{0, \dots, K-1\}$, let $\mathcal{C}_{I_i}(t) \subseteq \mathcal{P}(t)$ represent the specific pool of completed tasks that satisfy scope $i$'s compatibility criteria, and let $n_i$ represent the minimum sample support threshold required to activate that tier.

A task is automatically routed to the strictest possible scope whose historical pool satisfies its support threshold:

\begin{equation} 
\ell(t) = \min\big\{\, i \in \{0,\dots,K-1\} : |\mathcal{C}_{I_i}(t)| \ge n_i \,\big\} \label{eq:routing_logic} \end{equation}

The final baseline fallback tier, $I_{K-1}$, uses matching system history and requires one completed task. If no tier satisfies its support threshold, the task enters a cold-start state and receives an administrative bootstrap default reservation.

The selected scope, or the cold-start state, is recorded as the task's history\_scope and supplied as a model input feature. It does not partition the peak-memory models or restrict adjacent historical pools from contributing multi-scale statistical features to the broader history vector $\mathbf{h}_t$. The calibration layer evaluates the same hierarchy independently over residual history.

\subsection{Proxy Consumption Estimator}

\toolname{} trains a single consumption model $g$, global across scopes
but necessarily restricted to \emph{audited} tasks, since $c_\tau$ is
observed only when $q_\tau = 1$. 
Its training set is:

\begin{equation}
\{(\mathbf{s}_\tau, \mathbf{h}_{m,\tau}) \mapsto c_\tau : q_\tau = 1\}
\end{equation}

Note that $g$ conditions on the memory history $\mathbf{h}_{m,t}$ only,
matching the inputs available at both training and inference time.
Pooling audited tasks across all scopes gives $g$ substantially more
training data than any individual process. During development, its
out-of-fold prediction $\hat{c}_t$ is appended as an input feature to the
peak-memory models; during inference, the fitted global model supplies the
same feature. The predictive accuracy of $g$ therefore depends directly on
\emph{which} tasks the auditing policy selects, which we discuss in Section~\ref{sec:gates}.

\subsection{Peak-Memory Quantile Models}

CAMP fits a global Quantile LightGBM (QLGBM) ensemble for the point prediction and independent upper-tail ensembles for each workflow-process pair. Each process model is trained exclusively on rows belonging to that pair. During online learning, a workflow model provides a broader upper-tail fallback when available, followed by an independently trained global fallback. The regression target is the true post-execution peak memory footprint $y_t$, optimized under pinball loss across a vector of target quantiles $\mathcal{Q} = \{0.50, 0.90, 0.95, 0.99, 0.995\}$:

\begin{equation} f : \big(\mathbf{s}_t, \mathbf{h}_t, \hat{c}_t\big) \longmapsto \{Q_q\}_{q \in \mathcal{Q}} \label{eq:qlgbm_mapping} \end{equation}

\noindent where $Q_q$ represents the predicted $q$-quantile of peak memory. To minimize variance and stabilize predictions, each individual quantile target is modeled by an ensemble of three models initialized with distinct seeds, and the median value is selected. A sorting pass is executed immediately following inference to eliminate quantile crossing.

\subsection{Inference and Tail Calibration}
\label{sec:calibration}

During online execution, the global proxy model first appends $\hat{c}_t$. The global point model and the most specific available upper-tail model then output the ordered quantile set. A workflow process model is preferred, followed by the global fallbacks.

Because scientific workload memory footprints are highly heavy-tailed, \toolname{} considers candidate base requests from the point prediction, nearest-neighbour peak history, and upper-tail predictions. Each candidate is bounded below by the point prediction. On development tasks, \toolname{} selects the candidate that satisfies the global and per-workflow process coverage requirements with the smallest request-to-peak ratio. The selected candidate becomes the base request $b_t$.

\toolname{} then processes $b_t$ through a calibration layer $\mathcal{K}$ that adjusts the request based on the multiplicative scale-free log-residuals of past allocations:
\begin{equation} r_\tau = \log(y_\tau / b_\tau) \label{eq:log_residual} \end{equation}

A positive log-residual indicates a historical model under-prediction. Residual records are aggregated continuously per scope. Let $\widehat{Q}_{\alpha}(\cdot)$ denote the empirical quantile operator evaluated over a target residual set, where $\alpha$ is selected with the base request on development tasks. The calibration layer selects the strictest residual scope that satisfies its support threshold and blends its residual profile with a broader parent tier $p(i)$:

\begin{equation} \rho_t = \eta_{i}\, \widehat{Q}_{\alpha}\big(r_{i}\big) + \big(1-\eta_{i}\big)\, \widehat{Q}_{\alpha}\big(r_{p(i)}\big) \label{eq:blended_residual} \end{equation}

where $I_0$, $I_1$, and $I_2$ use $I_3$ as the parent tier, $I_3$ uses $I_4$, and $I_4$ and $I_5$ use $I_5$. The weight $\eta_i$ is a local residual confidence weight derived from scope sample support and log-residual variance; the $I_2$ weight also includes nearest-neighbour distance. The final calibrated memory reservation request $\hat{y}_t$ is scaled exponentially and rounded upward to the nearest whole MiB boundary:
\begin{equation} \hat{y}_t = \left\lceil b_t \exp(\max(0, \rho_t)) \right\rceil_{\mathrm{MiB}} \label{eq:final_rounded_allocation} \end{equation}


\subsection{History Construction: $\mathbf{h}_{m,t}$, $\mathbf{h}_{c,t}$}
\label{sec:history}

Both models above condition on summarized history rather than raw
task-level records. For every scope $i \in \{0,\dots,K-1\}$---not just
$i=\ell(t)$---\toolname{} computes a summary block, so a task routed to a
strict scope still sees broader, more stable statistics as context. Each
block summarizes the peak-RSS values (and, where audited, consumption
values) of scope $i$'s pool into medians, upper-tail percentiles, spreads,
and sample support. 
% ; blocks from all $K$ scopes are concatenated, together
% with the routed index $\ell(t)$, into the fixed-width vectors
% $\mathbf{h}_{m,t}$ and $\mathbf{h}_{c,t}$. The budget constraint makes
% consumption history systematically thinner than memory history---at most a
% fraction $B$ of tasks are ever audited---so empty or under-supported
% blocks are median-imputed and paired with explicit missingness flags,
% ensuring the models can distinguish genuine observations from
% placeholders. 
Exact field definitions, confidence metrics, and imputation
mechanics are detailed in the accompanying technical
report~\cite{camp2026techreport}.

\subsection{Selective Auditing}
\label{sec:gates}

Full eBPF auditing supplies the consumed-data signal that extends
$\mathbf{h}_{c,t}$, but auditing every task incurs the overhead measured
in Section~\ref{sec:results}. The auditing policy $\pi$ therefore decides,
before launch, which pending tasks should spend the budget $B$. An audit
cannot help the task that produces it i.e., its value begins only when the
record becomes history for later tasks.

Each pending task is scored from an audit-feature vector $\mathbf{x}_t$
drawn from the same prelaunch-legal information available to the
allocator: $\mathbf{s}_t$, $\mathbf{h}_t$, the active request, and the
model's quantiles. The task's own $y_t$, $c_t$, runtime, and outcome are
excluded; they appear only as labels when scoring models are trained on
completed tasks. Because score scales differ across workflows, the model
outputs and discovery signals are converted to within-workflow percentile
ranks $\operatorname{rank}_w(\cdot) \in (0,1]$ over pending tasks of the
same workflow, while history confidence enters directly.

The budget is dominated by a single question: \emph{is the active memory
request likely to be too small?} Tasks whose requests are about to fail
are valuable audit targets because they concentrate consumed-data
observations exactly where the predictor was weakest. \toolname{} estimates
this with two lightweight models trained
on completed tasks: a classifier for the underallocation probability
$p_t^r = \Pr(y_t > \hat y_t \mid \mathbf{x}_t)$, and, fitted only on
completed underallocated tasks, a regressor for the conditional deficit
$\hat\delta_t = \mathbb{E}[y_t - \hat y_t \mid y_t > \hat y_t,
\mathbf{x}_t]$. The \emph{risk score}
\begin{equation}
R_t = 0.70\,\operatorname{rank}_w(p_t^r)
    + 0.30\,\operatorname{rank}_w(p_t^r\,\hat\delta_t)
\label{eq:audit-risk}
\end{equation}
ranks a task highly when failure is likely, when its expected deficit is
large, or both. At each decision point, \toolname{} derives an exact audit
count from $B\,|U|$, divides it among workflows in proportion to pending
counts, and spends 80\% of each workflow's quota on its highest-$R_t$
tasks.

The remaining 20\% assigns 10\% to a \emph{heavy-information} score
selecting tasks whose consumed-data measurement is likely to improve future
predictions and splits the final 10\% between \emph{discovery} and random
draws from workflow-process groups. The random
component protects against unmonitored drift and keeps inclusion
probabilities non-degenerate; score definitions and quota mechanics
appear in the technical report~\cite{camp2026techreport}.

